\documentclass{article}


\usepackage{mib}

\usepackage[T2A]{fontenc}
\usepackage{color}
\usepackage[utf8]{inputenc} % allow utf-8 input
\usepackage[T1]{fontenc}    % use 8-bit T1 fonts
\usepackage{hyperref}       % hyperlinks
\usepackage{url}            % simple URL typesetting
\usepackage{booktabs}       % professional-quality tables
\usepackage{amsfonts}       % blackboard math symbols
\usepackage{nicefrac}       % compact symbols for 1/2, etc.
\usepackage{microtype}      % microtypography
\usepackage{lipsum}
\usepackage{graphicx}
\usepackage{color}
\usepackage{amsthm,amsmath,amssymb}
\usepackage[english, russian]{babel}
%\PassOptionsToPackage{silent}{fontspec}
%\graphicspath{ {./images/} }
\setcounter{page}{1}


\title{An initial-boundary value problem for a degenerate partial differential equation including tempered Caputo fractional operator}


\author{
 Karimov Erkinjon \\
  Department of Mathematics and Informatics\\
Fergana State University\\
Fergana, Uzbekistan\\
erkinjon@gmail.com\\
  %% examples of more authors
   \And
 Murolimova Nargiza\\
  Department of Mathematics and Informatics\\
Fergana State University\\
Fergana, Uzbekistan\\
f.nargiza97@gmail.com\\
  \And
Usmonov Doniyor\\
  Department of Mathematics and Informatics\\
Fergana State University\\
Fergana, Uzbekistan\\
dusmonov909@gmail.com\\
  %% \AND
  %% Coauthor \\
  %% Affiliation \\
  %% Address \\
  %% \texttt{email} \\
  %% \And
  %% Coauthor \\
  %% Affiliation \\
  %% Address \\
  %% \texttt{email} \\
  %% \And
  %% Coauthor \\
  %% Affiliation \\
  %% Address \\
  %% \texttt{email} \\
}

\begin{document}
\maketitle
\begin{center}
{\large \sc \bf Abstract}
\end{center}
\begin{abstract}
{ The present research is devoted to studying a unique solvability of an initial-boundary problem for a degenerate partial differential equation involving tempered Caputo fractional derivative in the time variable. First, we solve the Cauchy problem for an ordinary fractional differential equation with the tempered Caputo derivative representing its solution via the Kilbas-Saigo function. The solution of the main problem is represented in series form and using the explicit form of the corresponding Cauchy problem, certain properties of the Kilbas-Saigo function, and imposing certain conditions on given functions, we have proved the existence of a unique solution.} 

\end{abstract}
\label{firstpage}

{\bf Keywords:} %\keywords
{Initial-boundary problem; Caputo tempered fractional derivative; degenerate differential equation; Kilbas-Saigo function.}

{\bf MSC 2020:} 35R11


\lhead{{\sf Karimov E., Murolimova N., Usmonov D.~An initial-boundary value problem for a degenerate ...}}

{\bf \large 1. Introduction}

Text of the introduction.

For the cite references, please use \cite{Author3}.


{\bf \large 2. Main part}

For definitions, remarks, propositions, theorems, lemmas, and corollaries, please use the following formats:  

\textbf{Definition 2.1.} 
\rm{Text of the definition.}

\textbf{Remark 2.2.} 
\rm{Text of the remark.}

\textbf{Proposition 2.3.} 
\emph{Text of the proposition.}

\medskip
{\sf Proof. } The proof of the proposition. 

The proof of Proposition 2.3 is complete.  
\hfill $\Box$


\textbf{Theorem 2.4.} 
\emph{Text of the theorem.}

\medskip
{\sf Proof. } The proof of the theorem. 

The proof of Theorem 2.4 is complete.  
\hfill $\Box$ 

\textbf{Lemma 2.5.}
\emph{Text of the lemma.}

\medskip
{\sf Proof. } The proof of the lemma. 

Lemma 2.5 has been proved.  
\hfill $\Box$


\textbf{Corollary 2.6.}
\emph{Text of the corollary.}

\medskip
{\sf Proof. } The proof of the corollary. 

Corollary 2.6 has been proved.  
\hfill $\Box$

\medskip
For each group of objects, use its own numbers. (e.g. Theorem 1.1, Theorem 1.2,... Lemma 1.1, Lemma 1.2, ... Corollary 1.1, Corollary 1.2,...)

\textbf{Use double dollars for separate formulas:}
$$
u_{xx}-u_{t}=f(x,t), \quad (x,t) \in \Omega_1.
$$

\textbf{The numeration of formulas are as follows:}
\begin{equation}\label{1.1}
u(x,0)=f_1(x), \quad x\in\mathbb{R},
\end{equation}
where $f_1(x)$ is a given function.


\textbf{For citing formulas please use the following format:} 

Using the condition \eqref{1.1}, we get...


\medskip
\textbf{For tabular, please use the following format:}

\bigskip
\begin{center}

\begin{tabular}{|c|c|c|}
	\hline 
	Number of the type changing lines  & Number of gluing conditions & Remarks \\ 
	\hline 
	1& 2 & According to the \\ &&order of the equation \\ && the number of \\ &&gluing conditions change\\ 
	\hline 
	2& 4 & According to the \\ &&order of the equation \\ && the number of \\ &&gluing conditions change\\
	\hline 
\end{tabular} 
\end{center}


For the references use the following formats: 

\bigskip

\textbf{\Large References}

\begin{enumerate}

% For books use the following format:
\bibitem{Author1} Authors.Title of the book. Year of publication. City: Publisher.
% e.g.:
%\bibitem{Evans}  Evans L.C. Partial differential equations. 2010. New York: American Mathematical Society.

% For articles use the following format:

\bibitem{Author2} Authors. Title of the article. The name of the Journal. Year. Vol.5, Issue 1,
pp. 72-78.
%\bibitem{Karimov}  Karimov E.T. Fundamental solutions for a class of three-dimensional elliptic equations with singular coefficients. Applied Mathematics Letters. 2009. Vol.22, Issue 12, pp. 1828-1832.

% For the abstracts of the Conference use the following format:
\bibitem{Author3} Authors. Title of the abstract. The name of the Conference. City, Country. Year, pp. 1009-1015.

%\bibitem{Karimov_2} Karimov E.T. Boundary value problems for mixed type equations. Modern problems of Mathematical Physics. Tashkent, Uzbekistan. March 12-13,2020, pp. 1009-1015.

% For archive preprints use the following format: 
\bibitem{Author4} Authors. arXiv preprint. arXiv:1305.4992. Year. 

\end{enumerate}


\newpage

% Title in Uzbek. Please, use latin letters.
\begin{center}
	{\sc Temperlangan Kaputo kasr tartibli hosila ishtirok etgan xususiy hosilali tenglama uchun boshlang`ich-chegaraviy masala }\\
	% Author's name in Uzbek
	{\bf Karimov Erkinjon, Murolimova Nargiza, Usmonov Doniyor}\\
\end{center}
% Abstract in Uzbek
\medskip

{Ushbu tadqiqot vaqt bo`yicha temperlangan Kaputo kasr tartibli hosilasi ishtirok etgan xususiy hosilali tenglama uchun bir boshlang`ich-chegaraviy masalaning bir qiymatli yechilishiga bag`ishlangan. Dastlab temperlangan Kaputo kasr tartibli hosilasi ishtirok etgan oddiy differensial tenglama uchun Koshi masalasini yechamiz va yechimini Kilbas-Saygo funskiyasi orqali ifodalaymiz. So`ngra asosiy masala yechimini qatro ko`rinishida ifodalaymiz va Koshi masalasi yechimi, Kilbas-Saygo funksiyasi xossalaridan foydalanib, berilganlarga ma'lum shartlar asosida tadqiq etilayotgan masalaning bir qiymatli yechilishini isbotlaymiz.}\\

\medskip
%Keywords in Uzbek
{{\bf Kalit so`zlar:} Boshlang`ich-chegaraviy masala; temperlangan Kaputo kasr tartibli hosilasi; buziladigan differensial tenglama; Kilbas-Saygo funksiyasi.}\\

\bigskip

% Title in Russian. Please, use encoding UTF-8 
\begin{center}
{\sc Начально-краевая задача для вырождающегося уравнения в частных производных с темперированной дробной производной Капуто }\\
% Author's name in Russian
{\bf Каримов Эркинжон, Муролимова Наргиза, Усмонов Дониер}\\
\end{center}
% Abstract in Russian
\medskip

{Настоящее исследование посвящено однозначной разрешимости одной начально-краевой задачи для вырождающегося уравнения в частных производных с темперированной дробной производной Капуто по временной переменной. Сначала мы решаем задачу Коши для обыкновенного дифференциального уравнения с темперированной дробной производной Капуто представляя решение функцией Килбас-Сайго. Решение основной задачи представлено в виде ряда и пользуясь решением задачи Коши, а также свойствами функции Килбас-Сайго, при определенных условиях на заданные, доказано однозначная разрешимость исследуемой задачи.}\\

\medskip
%Keywords in Russian
{{\bf Ключевые слова:} Начально-краевая задача; темперированная дробная производная Капуто; вырождающегося дифференциальное уравнение; функция Килбас-Сайго.}\\

\bigskip

\textbf{Received: **/**/202*}\\

\textbf{Accepted: **/**/202*}\\


\bigskip
\textcolor{red}{\Large Cite this article}

Karimov E., Murolimova N. and Usmonov D.~ {An initial-boundary value problem for a degenerate partial differential equation including tempered Caputo fractional operator.}  \textit{Bull.~Inst.~Math.},~ 2023, Vol.6, No 1, pp. \pageref{firstpage}-\pageref{lastpage}


\label{lastpage}

\end{document}


